% TOP_PROBLEM: {"id": 30007294, "title": "Abstract Cuntz-semigroup criterion for Jiang–Su absorption", "category_id": 8, "category": {"id": 8, "name": "analysis", "display_name": "Analysis", "description": "Limits, continuity, calculus, and function theory.", "slug": "analysis", "order_index": 8, "created_at": "2026-07-31T15:26:25.671Z"}, "status": "open"}
% FIRST_POSTED: 2026-10-06
% ATTEMPT_STATUS: solved
% OPENAI_CLAIM: {"schema_version":1,"source_repo":"https://github.com/openai/math","source_commit":"adc7f1241b42e322a6451854ab7e4b4c146bf78a","release_date":"2026-10-06","match":"full","resolution":"proved","independently_reviewed":false,"second_pass":"agreed","families":[{"family":"291","title":"Cuntz comparison, nuclear dimension, and equivariant Jiang–Su stability","subject":"Operator algebras","manuscripts":[{"title":"Cuntz comparison and Jiang–Su absorption","pdf_path":"preprints/Cuntz-comparison-and-Jiang-Su-absorption-September-23-2026/paper.pdf","readme_path":"preprints/Cuntz-comparison-and-Jiang-Su-absorption-September-23-2026/README.md","date":"2026-09-23","theorem_ref":"ext:thm:main (sections/introduction.tex:61); thm:condition-C (sections/introduction.tex:127); cor:canonical (sections/introduction.tex:135); cor:abstract-cu (sections/introduction.tex:154)","statement_sources":[{"path":"preprints/Cuntz-comparison-and-Jiang-Su-absorption-September-23-2026/build/sections/introduction.tex","line":61,"label":"ext:thm:main"},{"path":"preprints/Cuntz-comparison-and-Jiang-Su-absorption-September-23-2026/build/sections/introduction.tex","line":127,"label":"thm:condition-C"},{"path":"preprints/Cuntz-comparison-and-Jiang-Su-absorption-September-23-2026/build/sections/introduction.tex","line":135,"label":"cor:canonical"},{"path":"preprints/Cuntz-comparison-and-Jiang-Su-absorption-September-23-2026/build/sections/introduction.tex","line":154,"label":"cor:abstract-cu"}]}],"lean":{"doc_path":"lean/docs/291.md","comparators":[],"declarations":[],"covers_main_theorem":false},"reasoning_trace":null}]}
\documentclass[11pt]{article}
\usepackage[margin=25mm]{geometry}
\usepackage{fontspec}
\setmainfont{Latin Modern Roman}
\usepackage{amsmath,amssymb,mathtools}
\usepackage{unicode-math}
\setmathfont{Latin Modern Math}
\usepackage{enumitem}
\usepackage{xurl,hyperref,bookmark}
\hypersetup{colorlinks=true,urlcolor=blue}
\setcounter{secnumdepth}{0}
\setlength{\parindent}{0pt}
\setlength{\parskip}{5pt}
\setlength{\emergencystretch}{3em}
\begin{document}
\section*{Abstract Cuntz-semigroup criterion for Jiang–Su absorption}

\subsection{OpenAI's claim}
A separable nuclear C*-algebra that is simple and non-elementary is claimed to absorb Jiang--Su under extended-functional strict comparison, including nonunital and unbounded-trace cases; semigroup divisibility also yields a nonsimple criterion. The cited manuscript argues for that result; its relation to this problem is specified below.

\subsection{Mathematical statement of the claimed result}
\paragraph{Mathematical claim 1.} Let $\operatorname{Cu}(A)$ be the ordered Cuntz semigroup of positive elements of $A\otimes\mathcal K$, modulo mutual norm-approximation comparison. For separable simple nuclear non-elementary $A$, OpenAI claims $A\simeq A\otimes_{\min}\mathcal Z$ if, for every $x,y\in\operatorname{Cu}(A)$, inequalities $\varphi(x)\le\varphi(y)$ for all additive order/supremum-preserving functionals $\varphi:\operatorname{Cu}(A)\to[0,\infty]$, strict whenever $0<\varphi(y)<\infty$, imply $x\le y$.
It further claims absorption for every nonzero separable nuclear $A$ whose Cuntz semigroup satisfies both
\[
(l+1)x\le ly\ \Longrightarrow\ x\le y\quad(l\ge1),\qquad
\forall x,\ N\ge1\ \exists u:\ Nu\le x\le(N+1)u.
\]
In that case a unital $\mathcal Z\to (A_\omega\cap A')/\operatorname{Ann}(A)$ is asserted for every free ultrafilter, where $A_\omega$ is the norm ultrapower and $\operatorname{Ann}(A)=\{z:zA=Az=0\}$. Finally it claims
\[
A\simeq A\otimes_{\min}\mathcal Z\quad\Longleftrightarrow\quad
\operatorname{Cu}(A)\simeq\operatorname{Cu}(A\otimes_{\min}\mathcal Z)
\]
for all separable nuclear $A$, allowing an abstract ordered-semigroup isomorphism rather than only the canonical first-factor map. Here $\mathcal Z$ is the Jiang--Su algebra, and non-elementary excludes compact-operator algebras.

\subsection{Outline of the argument}
Manuscript 1 (family 291, 2026-09-23): Trace transfer into small hereditary supports followed by spectral rounding establishes full Cuntz-semigroup divisibility. Pure-state compression produces norm-visible matrix cones, while localized nuclear approximations produce trace-visible ones. Comparison packs their defects into a central corner, constructing dimension-drop maps whose lifted intertwining relations give Jiang--Su absorption.

\subsection{Scope relative to this problem}
Manuscript 1 (family 291, 2026-09-23): The problem selects the mathematical statement summarized here, including its stated hypotheses and parameter range. The cited theorem gives a claimed proof of that assertion in full.

\subsection{Formal verification}
For family 291, the scope page documents the following selections. The selection gives uniform Gamma from real rank zero of the designated tracial ultrapower. The documented comparison scope also treats Jiang--Su absorption and Gamma under strict comparison, with additional trace and target-rank conditions in the projectionless case. Manuscript 1 is not among the papers selected by that page, so the family link does not formalize its main conclusion. It does not establish formal coverage of the complete problem clause.

This site has not executed the Lean files or checked their dependencies. These descriptions report the scope of the linked declarations, rather than a new verification of the proofs.

\subsection{Status caveat}
OpenAI attributes these results to an internal model. Its release notice cautions that results lacking full formalization can contain mistakes. This site has not independently assessed the proofs; a solved label records the scope of the reported claim.

\subsection{Sources}
\begin{itemize}
\item OpenAI, \emph{Cuntz comparison and Jiang–Su absorption}, 2026-09-23; citation key \texttt{OAI:Cuntz-comparison-and-Jiang-Su-absorption-September-23-2026}. \href{https://github.com/openai/math/blob/adc7f1241b42e322a6451854ab7e4b4c146bf78a/preprints/Cuntz-comparison-and-Jiang-Su-absorption-September-23-2026/paper.pdf}{PDF (GitHub)}; \href{https://github.com/openai/math/raw/adc7f1241b42e322a6451854ab7e4b4c146bf78a/preprints/Cuntz-comparison-and-Jiang-Su-absorption-September-23-2026/paper.pdf}{PDF (direct)}; \href{https://github.com/openai/math/blob/adc7f1241b42e322a6451854ab7e4b4c146bf78a/preprints/Cuntz-comparison-and-Jiang-Su-absorption-September-23-2026/README.md}{Manuscript page and citation}; \href{https://github.com/openai/math/blob/main/preprints/Cuntz-comparison-and-Jiang-Su-absorption-September-23-2026/paper.pdf}{Latest version}; Statement source: \href{https://github.com/openai/math/blob/adc7f1241b42e322a6451854ab7e4b4c146bf78a/preprints/Cuntz-comparison-and-Jiang-Su-absorption-September-23-2026/build/sections/introduction.tex\#L61}{ext:thm:main}; Statement source: \href{https://github.com/openai/math/blob/adc7f1241b42e322a6451854ab7e4b4c146bf78a/preprints/Cuntz-comparison-and-Jiang-Su-absorption-September-23-2026/build/sections/introduction.tex\#L127}{thm:condition-C}; Statement source: \href{https://github.com/openai/math/blob/adc7f1241b42e322a6451854ab7e4b4c146bf78a/preprints/Cuntz-comparison-and-Jiang-Su-absorption-September-23-2026/build/sections/introduction.tex\#L135}{cor:canonical}; Statement source: \href{https://github.com/openai/math/blob/adc7f1241b42e322a6451854ab7e4b4c146bf78a/preprints/Cuntz-comparison-and-Jiang-Su-absorption-September-23-2026/build/sections/introduction.tex\#L154}{cor:abstract-cu}.
\item Family 291: \href{https://github.com/openai/math/blob/adc7f1241b42e322a6451854ab7e4b4c146bf78a/lean/docs/291.md}{Lean scope}.
\item \href{https://github.com/openai/math/blob/adc7f1241b42e322a6451854ab7e4b4c146bf78a/CONTENTS.md}{OpenAI catalogue entry}.
\end{itemize}
\end{document}
